\section{Cosmological scales associated with the 108-state cycle and the Perron operator}
\label{rm:ch:cosmo}

\subsection{Typological separation of the \(2\) cosmological realizations}

The construction in cosmology comprises two internal quantities of
distinct provenance. The first is derived from the periodic CT108 closure and
determines a curvature scale. The second is derived from the Perron cocycle and
determines a law of memory dilution.
Both use the same trunk \(\APP\to\TRIT\to\TPK\), but do not have the
same domain or the same selector:

{\footnotesize
\begin{equation}
\TPK\longrightarrow\text{enriched state}
\longrightarrow
\begin{cases}
\text{chronological structure }108\longrightarrow
R_{108}\longrightarrow(H_{108},\Lambda_{108},T_{\rm dS},S_{\rm dS}),\\
\text{Perron cocycle}\longrightarrow
(a_n,M_n)\longrightarrow(H_\varphi,s^2\propto a^{-6}).
\end{cases}
\label{rm:eq:cosmo-two-branches}
\end{equation}
}

The eventual equality of units or order of magnitude does not warrant
identifying these realizations. The first falsification criterion of the
chapter is the attempt to identify them without an intertwiner that
preserves the chronological structure, orientation,
memory, and the dimensional chart.

\subsection{Periodicity condition of the 108-state cycle}

Let \(\ell_0=c_{\rm int}t_0\). The period and angular frequency of the return
are

\begin{equation}
T_{108}=108t_0,
\qquad
\omega_{108}=\frac{2\pi}{108t_0},
\qquad
R_{108}=\frac{c_{\rm int}}{\omega_{108}}
=\frac{54}{\pi}\ell_0.
\label{rm:eq:cosmo-R108}
\end{equation}

The length of the full turn is

\begin{equation}
C_{108}=2\pi R_{108}=108\ell_0.
\label{rm:eq:cosmo-C108}
\end{equation}

For either of the two oriented action sections
\(a\in\{\mathrm{pre},\mathrm{ret}\}\), the circular quantum

\[
E_{108,a}=\hbar_a\omega_{108},
\qquad
m_{108,a}=E_{108,a}/c_{\rm int}^2
\]

satisfies simultaneously

\begin{equation}
\bar\lambda_C(m_{108,a})=R_{108},
\qquad
\lambda_C(m_{108,a})=C_{108}.
\label{rm:eq:cosmo-compton}
\end{equation}

The circular realization therefore does not take an external Planck length.
It first constructs a radius and then demonstrates that the reduced and
complete Compton readings are its two realizations.

\subsection{Gravitational selector and Planck--TRIT relation}

The homogeneous gravitational transducer is

\begin{equation}
G_a(\theta)=\theta\frac{c_{\rm int}^5t_0^2}{\hbar_a}.
\label{rm:eq:cosmo-Gtheta}
\end{equation}

On the circular quantum, the reduced gravitational radius is equal to

\[
r_{{\rm g},a}(\theta)
=\frac{G_a(\theta)m_{108,a}}{c_{\rm int}^2}
=\theta\frac{\pi}{54}\ell_0.
\]

The condition \(r_{{\rm g},a}=R_{108}\) selects in a unique mode

\begin{equation}
\theta_{108}=\left(\frac{54}{\pi}\right)^2.
\label{rm:eq:cosmo-theta108}
\end{equation}

The coordinate does not depend on the action sheet: upon multiplying
\(G_a(\theta_{108})\hbar_a\), the pre/return orientation cancels. From here
there result

\begin{align}
\ell_P&=R_{108}=\frac{54}{\pi}\ell_0,
&
t_P&=\omega_{108}^{-1}=\frac{54}{\pi}t_0,
\label{rm:eq:cosmo-planck-length-time}\\
E_P&=\hbar_a\omega_{108}
=\frac{\pi}{54}\frac{\hbar_a}{t_0},
&
E_P&=k_B\Theta_{\rm clk}\log3.
\label{rm:eq:cosmo-planck-trit}
\end{align}

The \cref{rm:eq:cosmo-planck-trit} expresses the structural confluence:
\(\log3\) is the informational content of a TRIT decision and, in the
thermal chart, converts chronological temperature into the Planck
quantum. The identity does not equate the TRIT state space with gravitational space;
it proves that two functional realizations determine the same scalar through distinct
maps.

\subsection{Conditioned realization of de Sitter geometry}

Suppose now that the Palatini--Cartan--Holst realizer selects the internal
radius \(R_{108}\) as the radius of a four-dimensional de Sitter background. The
hypothesis is geometric and remains explicit: it is not deduced from a numerical
coincidence with \(H_0\). Under this hypothesis,

\begin{equation}
H_{108}=\frac{c_{\rm int}}{R_{108}}
=\frac{\pi}{54t_0},
\qquad
\Lambda_{108}=\frac{3}{R_{108}^2}
=\frac{\pi^2}{972\ell_0^2}.
\label{rm:eq:cosmo-ds-H-Lambda}
\end{equation}

The Gibbons--Hawking temperature and horizon entropy are

\begin{align}
k_BT_{\rm dS}
&=\frac{\hbar_a H_{108}}{2\pi}
=\frac{\hbar_a}{108t_0},
&
\frac{T_{\rm dS}}{\Theta_{\rm clk}}
&=\frac{\log3}{2\pi},
\label{rm:eq:cosmo-ds-temperature}\\
\frac{S_{\rm dS}}{k_B}
&=\frac{\pi R_{108}^2}{\ell_P^2}=\pi.
\label{rm:eq:cosmo-ds-entropy}
\end{align}

The \(\pi\) of \cref{rm:eq:cosmo-ds-entropy} is not introduced as the value of an
entropy to be reproduced: it appears because the circular selector gives
\(R_{108}=\ell_P\), and the area of the horizon sphere contains the
geometric factor \(4\pi\).

\begin{proposition}[Dimensionless family CT108--de Sitter]
Under the declared de Sitter realization, the quadruple
\[
\left(
H_{108}t_0,
\Lambda_{108}\ell_0^2,
\frac{T_{\rm dS}}{\Theta_{\rm clk}},
\frac{S_{\rm dS}}{k_B}
\right)
=\left(
\frac\pi{54},
\frac{\pi^2}{972},
\frac{\log3}{2\pi},
\pi
\right)
\]
is a dimensionless exact output.
\end{proposition}

\begin{falsifier}
The de Sitter realization is invalidated if the selected background does not
satisfy \(R_{\mu\nu}=\Lambda_{108}g_{\mu\nu}\), if the PCH realizer determines
a radius other than \(R_{108}\), or if the charts used in
\cref{rm:eq:cosmo-ds-H-Lambda,rm:eq:cosmo-ds-temperature} do not share
\((c_{\rm int},t_0,\hbar_a)\). None of these failures alters the
circular theorem \cref{rm:eq:cosmo-R108,rm:eq:cosmo-theta108}.
\end{falsifier}
